\exercisesection{§ 1 的习题}

下文考虑的所有根系都相对于实向量空间。我们用 \((x|y)\) 表示在 Weyl 群下不变的标量积（参见第 3 号）。

\begin{enumerate}
\def\labelenumi{\arabic{enumi})}
\item
  令 \(\mathbb{R}\) 为根系，\(\mathbb{R} = \mathbb{R}_1 \cup \mathbb{R}_2\) 为 \(\mathbb{R}\) 的一个划分。假设，当 \(x, y\) 是 \(\mathbb{R}_i\) 的两个元素且 \(x + y\)（分别为 \(x - y\)）是一个根时，就有 \(x + y \in \mathbb{R}_i\)（分别为 \(x - y \in \mathbb{R}_i\)），其中 \(i = 1, 2\)。那么，\(\mathbb{R}\) 是 \(\mathbb{R}_1\) 与 \(\mathbb{R}_2\) 的直和。（利用第 3 号定理 1 的推论，证明若 \(x \in \mathbb{R}_1\) 且 \(y \in \mathbb{R}_2\)，则 \((x|y) = 0\)。）
\item
  令 R 为根系，\(\alpha\) 和 \(\beta\) 为两个根。若 t 是一个标量且 \(\beta + t\alpha \in R\) ，则 \(2t \in Z\) 。（因为 \(n(\beta + t\alpha, \alpha) = n(\beta, \alpha) + 2t\) 。）若 \(\alpha\) 不可整除，则 \(t \in Z\) 。（否则利用命题 9 证明存在一个根 \(\gamma\)，它与 \(\alpha\) 正交，使得 \(\gamma + \frac{1}{2}\alpha \in R\)：于是 \((\gamma + \frac{1}{2}\alpha|\alpha) > 0\) ，所以 \(\frac{1}{2}\alpha \in R\)。）
\item
  令 R 为 V 中的不可约非约化根系。V 上存在一个非退化对称双线性型 \(\Phi\)，使得当将 V 与 \(V^{*}\) 通过 \(\Phi\) 识别时，有 \(R = R^{-}\)。（利用命题 13。）
\item
  令 \(\Gamma\) 为有限秩自由 \(\mathbf{Z}\)-模，其秩为 \(l\)，\(\Gamma^{*}\) 为其对偶 \(\mathbf{Z}\)-模，I 为有限指标集，\((x_{i}, x_{i}^{*})_{i \in \mathbf{I}}\) 为 \(\Gamma \times \Gamma^{*}\) 中的元素族，且满足 \(\langle x_{i}, x_{i}^{*} \rangle = 2\) 对所有 \(i \in \mathbf{I}\) 成立，\(s_{i} = s_{x_{i}, x_{i}^{*}}\)。令 \(\mathbf{V}\) 为实向量空间 \(\Gamma \otimes_{\mathbf{Z}} \mathbf{R}\)，其对偶 \(\mathbf{V}^{*}\) 可与 \(\Gamma^{*} \otimes_{\mathbf{Z}} \mathbf{R}\) 识别。我们也用 \(s_{i}\) 表示反射 \(s_{i} \otimes 1\) 在 \(\mathbf{V}\) 中的作用。令 \(\mathbf{R}\)（分别为 \(\mathbf{R}'\)）为 \(x_{i}\)（分别为 \(x_{i}^{*}\)）的集合。令 \(\mathbf{E}\)（分别为 \(\mathbf{E}'\)）为 \(\mathbf{V}\)（分别为 \(\mathbf{V}^{*}\)）中由 \(\mathbf{R}\)（分别为 \(\mathbf{R}'\)）生成的子空间。假设 \(s_{i}(\mathbf{R}) = \mathbf{R}\) 对所有 \(i\) 成立。那么，\(\mathbf{R}\) 是 \(\mathbf{E}\) 中的根系。\(\mathbf{E}'\) 可自然地识别为 \(\mathbf{E}\) 的对偶，且 \(x_{i}^{*}\) 可与 \(x_{i}^{-}\) 识别，对所有 \(i\) 均如此。若 \(\mathbf{F}\)（分别为 \(\mathbf{F}'\)）是 \(\mathbf{V}\)（分别为 \(\mathbf{V}^{*}\)）中在所有 \(s_{i}\)（分别为 \(^t s_{i}\)）下不变的点所成的子空间，则 \(\Gamma \cap \mathbf{F}\)（分别为 \(\Gamma^{*} \cap \mathbf{F}'\)）生成 \(\mathbf{F}\)（分别为 \(\mathbf{F}'\)）。若置 \(\Gamma_{1} = (\Gamma \cap \mathbf{E}) \oplus (\Gamma \cap \mathbf{F})\)（分别为 \(\Gamma_{1}^{*} = (\Gamma^{*} \cap \mathbf{E}') \oplus (\Gamma^{*} \cap \mathbf{F}')\)），则 \(\Gamma / \Gamma_{1}\) 与 \(\Gamma^{*} / \Gamma_{1}^{*}\) 是同构的有限群。
\item
  令 \(R\) 为不可约根系。对任意 \(\beta \in R\)，
\end{enumerate}

\[
16 \gamma (\mathrm{R}) = (\sum_ {\alpha \in \mathrm{R}} n (\alpha , \beta) ^ {2}) (\sum_ {\alpha \in \mathrm{R}} n (\beta , \alpha) ^ {2})
\]

从而 \(\gamma (\mathbb{R}) = \gamma (\lambda R)\) 对任意非零标量 \(\lambda\) 成立。（利用第 12 号的公式 (17) 和 (19)。）

\begin{enumerate}
\def\labelenumi{\arabic{enumi})}
\setcounter{enumi}{5}
\tightlist
\item
  \begin{enumerate}
  \def\labelenumii{\alph{enumii})}
  \tightlist
  \item
    令 A 为有限型阿贝尔群，T 为 A 的一个不含 0 的有限子集。存在 A 的有限指数子群 H，使得 \(H \cap T = \varnothing\)。（令 \(t \in T\)。利用有限型阿贝尔群的结构，构造 A 的一个不含 t 的有限指数子群。然后对 \(\text{Card}(T)\) 作归纳。）
  \end{enumerate}
\end{enumerate}

\begin{enumerate}
\def\labelenumi{\alph{enumi})}
\setcounter{enumi}{1}
\tightlist
\item
  令 R 为根系，P 为 R 的一个闭对称子集。存在 Q(R) 的有限指数子群 H，使得 \(P = H \cap R\)。（对 Q(R) 中由 P 生成的子群取商，并利用 a) 和命题 23。）
\end{enumerate}

\begin{enumerate}
\def\labelenumi{\arabic{enumi})}
\setcounter{enumi}{6}
\item
  根系的连通指标等于其 Cartan 矩阵的行列式。（利用第 10 号的公式 (14)。另一方面证明：在带有标量积的实向量空间中，若 \((\varepsilon_1, \ldots, \varepsilon_n), (\varepsilon_1', \ldots, \varepsilon_n')\) 是满足 \((\varepsilon_i | \varepsilon_j') = \delta_{ij}\) 的两个基，则 \(\det_{(\varepsilon_1, \ldots, \varepsilon_n)} (\varepsilon_1', \ldots, \varepsilon_n') > 0\)。）
\item
  令 R 为约化根系，C 为一个室，\(2\sigma\) 为 \(R^{-}\) 的正元素之和，\(P'(R)\) 为满足 \(x \in P(R)\) 且 \(\langle x, \sigma \rangle \in Z\) 的集合。于是，\(Q(R) \subseteq P'(R) \subseteq P(R)\)，且 \(P(R)/P'(R)\) 的阶为 1 或 2。令 \((\beta_{1}, \ldots, \beta_{l})\) 为与 C 对应的 \(R^{-}\) 的基，并置 \(2\sigma = n_{1}\beta_{1} + \cdots + n_{l}\beta_{l}\)，其中 \(n_{i}\) 是 \textgreater0 的整数。则 \(P(R) = P'(R)\) 当且仅当所有 \(n_{i}\) 都是偶数。
\item
  令 \(\mathbb{R}\) 为根系，\(\mathbb{C}\) 为一个室。置 \(\mathrm{B}(\mathbb{C}) = \{\alpha_1, \ldots, \alpha_l\}\)，
\end{enumerate}

\[
\mathrm{R} _ {+} (\mathrm{C}) = \left\{\alpha_ {1}, \dots , \alpha_ {l}, \alpha_ {l + 1}, \dots , \alpha_ {s} \right\}
\]

（其中 \(\alpha_{i}\) 两两不同），并令 \(\alpha = \alpha_{1} + \dots +\alpha_{s}\)。

\begin{enumerate}
\def\labelenumi{\alph{enumi})}
\tightlist
\item
  对 \(i = 1, 2, \ldots, s\) ，置 \(\varepsilon_{i} = \pm 1\)。令 \(\alpha^{*} = \varepsilon_{1}\alpha_{1} + \cdots + \varepsilon_{s}\alpha_{s}\)。若 \((\alpha^{*}|\alpha_{i}) > 0\) 对 \(i = 1, \ldots, l\) 成立，则 \(\alpha^{*} = \alpha\) 且对所有 i 都有 \(\varepsilon_{i} = 1\)。（令 \(\gamma\)（分别为 \(\delta\)）为 \(\alpha_{i}\) 之和，其中 \(\varepsilon_{i} = 1\)（分别为 -1）。则 \((\gamma|\alpha_{i}) - (\delta|\alpha_{i}) > 0\)，\((\gamma|\alpha_{i}) + (\delta|\alpha_{i}) = (\alpha_{i}|\alpha_{i})\)，所以
\end{enumerate}

\(2(\gamma|\alpha_{i})>(\alpha_{i}|\alpha_{i});\) 由于 \(2(\gamma|\alpha_{i}) / (\alpha_{i}|\alpha_{i})\in \mathbf{Z},(\gamma |\alpha_{i})\geqslant (\alpha_{i}|\alpha_{i}) = (\alpha |\alpha_{i}),\) 因而 \(\gamma -\alpha \in \overline{\mathrm{C}}. \gamma \geqslant \alpha ,\) \(\gamma = \alpha\) 且 \(\delta = 0.)\)

\begin{enumerate}
\def\labelenumi{\alph{enumi})}
\setcounter{enumi}{1}
\tightlist
\item
  对 \(i = 1,2,\ldots,s\)，令 \(\varepsilon_{i} = \pm 1\)。\(\varepsilon_{i}\alpha_{i}\)（\(>0\)）所成的集合是相对于某个室的正根集合，当且仅当 \(\alpha^{*} = \sum_{i}\varepsilon_{i}\alpha_{i}\) 属于某个室。（为证明充分性，令 \(\mu_{i} = \pm 1\) 满足 \((\alpha^{*}|\mu_{i}\alpha_{i}) > 0\)。存在 \(w\in W(R)\)，将 \(\mu_{i}\alpha_{i}\) 的集合变为 \(\alpha_{i}\) 的集合，从而将 \(\alpha^{*}\) 变为 \(\alpha^{**} = \sum_{i}\varepsilon_{i}\mu_{i}\alpha_{\sigma(i)}\)，其中 \(\sigma \in \mathfrak{S}_s\)。应用 \(a\)，证明 \(\alpha^{**} = \alpha\)，所以 \(\mu_{i}\varepsilon_{i} = 1\)。)
\end{enumerate}

\begin{enumerate}
\def\labelenumi{\arabic{enumi})}
\setcounter{enumi}{9}
\item
  令 \(\alpha\) 为不可约约化根系 R 相对于某个基的最高根。当且仅当所有根长度相同时，\(\alpha^{-}\) 才是 \(R^{-}\) 的最高根。
\item
  采用第 11 号命题 33 的记号，证明 \(\theta_{i} + \theta_{j}\) 对任意一对 \((i,j)\) 都不是根。
\item
  令 R 为 V 中秩 \(\geqslant 3\) 的根系，\((\alpha_{1},\ldots,\alpha_{l})\) 为 R 的一个基。令 \(V'\)（分别为 \(V''\)）为 V 中由 \(\alpha_{i}\) 生成的子空间，其中 \(i \geqslant 2\)（分别为 \(i \geqslant 3\)）。令 \(R' = R \cap V'\)，\(R'' = R \cap V''\)。令 d（分别为 \(d', d''\)）为 R（分别为 \(R', R''\)）的 Cartan 矩阵的行列式。假设 \(\alpha_{1}\) 与所有 \(\alpha_{i}\)（除 \(\alpha_{2}\) 外）正交，且 \(\| \alpha_{1} \| = \| \alpha_{2} \|\)。证明 \(d = 2d' - d''\)。
\item
  令 R 为约化根系。\((\alpha_{1},\ldots,\alpha_{l})\) 为 R 的一个基，且 \(\alpha=c_{1}\alpha_{1}+\cdots+c_{l}\alpha_{l}\) 为一个根。则对所有 i 都有 \(c_{i}(\alpha_{i}|\alpha_{i})/(\alpha|\alpha)\in\mathbf{Z}\)。（考虑逆系统。）
\item
  令 R 为不可约根系，\(\lambda\) 为最大的根长。令 S 为由长度为 \(\lambda\) 且两两正交的根组成的 R 的子集所成的集合。则 S 的任意两个极大元素都可由 W(R) 的元素彼此变换。（利用第 V 章 §3 第 3 号的命题 11 和命题 1。）
\item
  令 \(R\) 为秩为 \(l\) 的根系。
\end{enumerate}

\begin{enumerate}
\def\labelenumi{\alph{enumi})}
\item
  当且仅当 R 含有 l 个两两强正交的根时，\(-1 \in \mathrm{W}(\mathrm{R})\)。（为证明必要性，利用第 V 章 §3 第 3 号的命题 1 对 l 作归纳。）
\item
  若 \(w \in \mathrm{W}(\mathrm{R})\) 的阶为 2，则存在一个由两两强正交的根组成的集合 S，使得 w 是 \(s_{\alpha} (\alpha \in \mathrm{S})\) 的乘积。
\end{enumerate}

\begin{enumerate}
\def\labelenumi{\arabic{enumi})}
\setcounter{enumi}{15}
\tightlist
\item
  令 \(\mathbb{R}\) 为根系。令 B 为 \(\mathbb{R}\) 的一个基，\(\mathbb{R}_+\) 为相对于此基的正根集合。对 \(w \in W(\mathbb{R})\)，置 \(F_w = R_+ \cap w(-R_+)\)。证明映射 \(w \mapsto F_w\) 是从 \(W(\mathbb{R})\) 到以 \(F\) 为 \(R_+\) 的子集、且 \(F\) 与 \(R_+ - F\) 都闭的集合的双射。（将命题 20 的推论 1 应用于子集 \(P = (R_+ - F) \cup (-F)\)。）
\end{enumerate}

¶ 17) 令 R 为约化根系，B 为 R 的一个基。对 B 的任意子集 J，令 W \(_{J}\) 为由反射 s \(_{\alpha}\)（其中 \(\alpha \in J\)）生成的 W(R) 的子群，并令 w \(_{J}\) 为 W \(_{J}\) 的最长元素。令 \(\Phi\) 为所有 w \(_{J}\) 所成的集合。

\begin{enumerate}
\def\labelenumi{\alph{enumi})}
\tightlist
\item
  证明，元素 \(w \in \mathrm{W}(\mathbb{R})\) 属于 \(\Phi\)，当且仅当
\end{enumerate}

\[
\mathrm{W} (\mathrm{B}) \subseteq \mathrm{R} _ {+} (\mathrm{B}) \cup (- \mathrm{B}).
\]

(为证明条件的充分性，令 \(\mathrm{J}_1\) 为 \(\alpha \in \mathrm{B}\) 且满足 \(w(\alpha) \in -\mathrm{B}\) 的元素所成的集合，并置 \(\mathrm{J} = -w(\mathrm{J}_1)\)。若 \(\alpha \in \mathrm{J}_1\)，则 \(w(\alpha) \in -\mathrm{J}\) 且 \(w_{\mathrm{J}}w(\alpha) \in \mathrm{B}\)；

若 \(\alpha \in \mathrm{B} - \mathrm{J}_1\)，则 \(w_{\mathrm{J}}w(\alpha) \in \mathrm{R}_{+}\)：否则 \(w(\alpha)\) 将为正的而 \(w_{\mathrm{J}}w(\alpha)\) 为负的，这将意味着 \(w(\alpha)\) 属于由 \(\beta \in J\) 生成的子系统，而 \(\alpha\) 属于由 \(\beta \in \mathrm{J}_1\) 生成的子系统，这是荒谬的。推出 \(w_{\mathrm{J}}w = 1\)。

\begin{enumerate}
\def\labelenumi{\arabic{enumi})}
\setcounter{enumi}{17}
\item
  令 R 为根系，P 为 R 的一个抛物子集。证明 P 在 R 中的补集是闭的。
\item
  令 \(\mathbb{R}\) 为根系，令 \(x\) 为 \(\mathrm{Q}(\mathbb{R})\) 中长度最小的非零元素。证明 \(x \in \mathbb{R}\)。
\end{enumerate}

¶ 20) 令 \((\alpha_{1},\ldots,\alpha_{l})\) 为不可约约化根系 R 的一个基。令 r 和 p 为两个整数，其中 \(p\geqslant2\)，并满足：

\[
\left(\alpha_ {1} \mid \alpha_ {1}\right) = \dots = \left(\alpha_ {r} \mid \alpha_ {r}\right) = p. \left(\alpha_ {r + 1} \mid \alpha_ {r + 1}\right) = \dots = p. \left(\alpha_ {l} \mid \alpha_ {l}\right).
\]

\begin{enumerate}
\def\labelenumi{\alph{enumi})}
\item
  令 \(\alpha = c_{1}\alpha_{1} + \cdots + c_{l}\alpha_{l}\) 为一个根。证明 \((\alpha|\alpha) = (\alpha_{1}|\alpha_{1})\)（换言之，\(\alpha\) 是长根）当且仅当 p 整除 \(c_{r+1},\ldots,c_{l}\)。
\item
  令 h 为 W(R) 的 Coxeter 数。证明 R 的最长根（分别为最短根）的数目等于 hr（分别为 \(h(l - r)\)）。
\end{enumerate}

\begin{enumerate}
\def\labelenumi{\arabic{enumi})}
\setcounter{enumi}{20}
\tightlist
\item
  令 R 为根系，B 为 R 的一个基。令 \(\alpha, \beta \in \mathrm{B}\) 和 \(w \in \mathrm{W}(\mathrm{R})\) 满足 \(\beta = w(\alpha)\)。证明存在一个由 R 的元素组成的序列 \(\alpha_1, \ldots, \alpha_n\)，以及一个序列 \(w_1, \ldots, w_{n-1}\)，其元素均属于 \(\mathrm{W}(\mathrm{R})\)，使得
\end{enumerate}

\begin{enumerate}
\def\labelenumi{(\roman{enumi})}
\item
  \(\alpha_{1} = \alpha, \quad \alpha_{n} = \beta.\)
\item
  \(w = w_{n - 1}\dots w_1\)
\item
  \(w_{i}(\alpha_{i}) = \alpha_{i + 1}\) 对 \(1\leqslant i\leqslant n - 1\) 成立
\item
  对所有 \(i \in [1, n - 1]\)，存在 \(\beta_i \in \mathbf{B}\)，使得 \(w_i\) 属于 \(\mathrm{W}(\mathbb{R})\) 中由反射 \(s_{\alpha_i}\) 和 \(s_{\beta_i}\) 生成的子群。
\end{enumerate}

¶ 22) 在第 11 号命题 33 的假设和记号下，以 \(\omega_{1},\ldots,\omega_{l}\) 表示 R 的基本权。

\begin{enumerate}
\def\labelenumi{\alph{enumi})}
\item
  证明 \(c^{-1}(\omega_i) = \omega_i - \theta_i\)。推得 \(1 - c\) 将 \(P(R)\) 映到 \(Q(R)\)。
\item
  令 \(f\) 为 \(\mathbb{R}\) 的连通指标，令 \(m_1, \ldots, m_l\) 为 \(W(\mathbb{R})\) 的指数。证明
\end{enumerate}

\[
f = \det (1 - c) = \prod_ {j = 1} ^ {l} \left(1 - \omega ^ {m _ {j}}\right) = 2 ^ {l} \prod_ {j = 1} ^ {l} \sin \left(m _ {j} \pi / h\right)
\]

其中 \(\omega = e^{2\pi i/h}\)。

\begin{enumerate}
\def\labelenumi{\alph{enumi})}
\setcounter{enumi}{2}
\tightlist
\item
  令 p 为素数。将 h 写成 \(h = p^{a}H\) 的形式，其中 H 不被 p 整除。~证明 p 整除 f，当且仅当 H 整除某个 \(m_{j}\)。
\end{enumerate}

\begin{enumerate}
\def\labelenumi{\arabic{enumi})}
\setcounter{enumi}{22}
\tightlist
\item
  令 R 为约化根系，X 为 P(R) 的一个子集。若 X 满足下述条件，则称 X 是饱和的：
\end{enumerate}

\begin{enumerate}
\def\labelenumi{(\Alph{enumi})}
\setcounter{enumi}{18}
\tightlist
\item
  对所有 \(p \in \mathbf{X}, \alpha \in \mathbb{R}\) 以及 \(i \in \mathbf{Z}\)，只要 \(i\) 介于 0 与 \(\langle p, \alpha \rangle\) 之间，就有 \(p - i\alpha \in \mathbf{X}\)。
\end{enumerate}

\begin{enumerate}
\def\labelenumi{\alph{enumi})}
\item
  证明每个饱和子集在 R 的 Weyl 群 W 下稳定。
\item
  证明，对 P(R) 的任意子集 A，都存在包含 A 的 P(R) 的最小饱和子集 S(A)。
\item
  令 C 为 R 的一个室，令 \(p \in \mathrm{P}(\mathrm{R}) \cap \overline{\mathrm{C}}\) 为支配权。令 \(S(p)\) 为包含 p 的 \(\mathrm{P}(\mathrm{R})\) 的最小饱和子集。~令 \(\Sigma(p)\) 为满足下述条件的元素 \(p' \in \mathrm{P}(\mathrm{R})\) 的集合：
\end{enumerate}

\begin{enumerate}
\def\labelenumi{(\roman{enumi})}
\item
  \(p \equiv p' \mod Q(R)\) .
\item
  对所有 \(w \in \mathbf{W}\)，\(w(p') \leqslant w\)（相对于由 C 定义的序关系）。
\end{enumerate}

证明 \(\Sigma(p)\) 是有限的，包含 p，并且是饱和的。由此推出 \(S(p)\) 包含于 \(\Sigma(p)\)。

\begin{enumerate}
\def\labelenumi{\alph{enumi})}
\setcounter{enumi}{3}
\tightlist
\item
  证明，若 \(\alpha\) 是 R 的最长元素，则 \(S(\alpha) = R \cup \{0\}\)。
\end{enumerate}

\begin{enumerate}
\def\labelenumi{\arabic{enumi})}
\setcounter{enumi}{23}
\tightlist
\item
  保持前一习题的记号和假设。
\end{enumerate}

\begin{enumerate}
\def\labelenumi{\alph{enumi})}
\tightlist
\item
  令 \(p \in \mathrm{P}(\mathbb{R})\)，令 \(\mathbf{W}.p\) 为 \(p\) 在 \(\mathbf{W}\) 下的轨道。证明下列条件等价：
\end{enumerate}

\begin{enumerate}
\def\labelenumi{(\roman{enumi})}
\item
  \(S(p) = W.p.\)
\item
  \(\langle p, \alpha^{\prime} \rangle = 0, 1\) 或 \(-1\) 对所有 \(\alpha \in \mathbb{R}\) 成立。
\end{enumerate}

（若 (ii) 不满足，构造一个元素 \(q \in S(p)\)，使得 \((q|q) < (p|p)\)，从而 \(q \notin W.p.\)）

\begin{enumerate}
\def\labelenumi{\alph{enumi})}
\setcounter{enumi}{1}
\item
  令 X 为 P(R) 的非空饱和子集。证明 X 包含一个满足上述条件 (i) 和 (ii) 的元素 p。（取 X 中长度最小的元素 p。）
\item
  假设 \(\mathbb{R}\) 不可约。令 \(\mathbb{C}\) 为 \(\mathbb{R}\) 的一个室，令 \(\mathbb{B} = \{\alpha_i\}_{i\in I}\) 为相应的基，并令
\end{enumerate}

\[
\gamma^ {*} = \sum_ {i} n _ {i} \alpha_ {i}
\]

为 \(R'\) 的最高根。令 J 为 \(i \in I\) 且满足 \(n_{i} = 1\) 的集合。令 \(p \neq 0\) 为支配权。证明 a) 的条件 (i) 和 (ii) 分别等价于下列每个条件：

\begin{enumerate}
\def\labelenumi{(\roman{enumi})}
\setcounter{enumi}{2}
\item
  \(\langle p, \gamma^{*} \rangle = 1\) .
\item
  存在 \(i \in \mathrm{J}\)，使得 \(p\) 等于相应的基本权。
\end{enumerate}

满足这些条件的权称为极小权。证明 \(P(R)\) 的每个非空饱和子集都包含 0 或一个极小权。

